\chapter{Local cohomology and its calculation}
\label{ch:local-cohomology}

Throughout this chapter \(R\) is Noetherian, \(I\) is an ideal of \(R\),
and modules are arbitrary unless a finiteness condition is given.
The module-localization formulas and exactness used below are those of
\cite[ALG-05-C01-S01]{p103}. We use the same sign for the connecting
homomorphism as in the cochain long exact sequence.

\section{The torsion functor and its derived functors}

\begin{definition}[Sections supported in an ideal]\label{res:L01}
For an \(R\)-module \(M\) put
\[
 \Gamma_I(M)=\bigcup_{n\geq 1}(0:_M I^n).
\]
This is an additive left exact functor of \(M\), depends only on \(\sqrt I\),
and equals the identity for \(I=0\) and zero for \(I=R\).
\end{definition}
\begin{proof}
The annihilator submodules form an ascending union. A map preserves
annihilation by \(I^n\), and on a submodule the union is its intersection
with \(\Gamma_I(M)\). These observations give functoriality and left
exactness. If an ideal \(J\) generated by \(g_1,\ldots,g_t\) lies in
\(\sqrt I\), choose \(b_j\geq1\) with \(g_j^{b_j}\in I\). Every product
of \(1+\sum_j(b_j-1)\) generators contains such a power, so some power
of \(J\) lies in \(I\). Apply this argument in both directions if the
radicals agree. The two extreme ideals give the stated functors directly.
\end{proof}

\begin{theorem}[Derived torsion and the Ext limit]\label{res:L02}
Let \(H_I^q=R^q\Gamma_I\) for \(q\geq0\), and put \(H_I^q=0\) for
negative \(q\). A short exact sequence of modules gives a natural long
exact sequence. There is a natural isomorphism
\[
 H_I^q(M)\simeq\varinjlim_{n}\operatorname{Ext}_R^q(R/I^n,M).
\]
Each \(H_I^q(M)\) is \(I\)-power torsion.
\end{theorem}
\begin{proof}
Choose an injective resolution \(M\to Q^\bullet\) using
\cite[ALG-03-C02-S03]{p103}. Apply the injective horseshoe and the
connecting morphism for cochain complexes to obtain the long exact
sequence. The natural maps \(R/I^{n+1}\to R/I^n\) yield compatible
inclusions
\[
 \operatorname{Hom}_R(R/I^n,Q^\bullet)
   \longrightarrow \operatorname{Hom}_R(R/I^{n+1},Q^\bullet).
\]
Evaluation at \(1\) identifies their union with \(\Gamma_I(Q^\bullet)\).
Filtered colimits of modules are exact because a finite relation between
representatives holds in the limit precisely when it holds at a sufficiently
late stage. Applied to cycles and boundaries this commutes cohomology
with the colimit. The complex at stage \(n\) is killed by \(I^n\) by
\(R\)-linearity, so its cohomology and the resulting union are power
torsion. Resolution comparison gives naturality, including the maps
for short exact sequences.
\end{proof}

\begin{lemma}[Torsion of an injective module]\label{res:L03}
If \(Q\) is injective, then \(\Gamma_I(Q)\) is injective.
\end{lemma}
\begin{proof}
Use Baer's criterion \cite[ALG-03-C02-R03]{p103}. For a map
\(u:J\to\Gamma_I(Q)\) from an ideal \(J\), its image is killed by one
power \(I^n\) because \(J\) is finite. Artin--Rees
\cite[ALG-05-C04-R02]{p103} gives \(N\) such that
\(J\cap I^N\subset I^nJ\). Consequently \(u\) extends by zero to
\(J+I^N\). Extend this new map to \(R\) with values in \(Q\).
Its value at \(1\) is killed by \(I^N\), so the entire extension lands
in \(\Gamma_I(Q)\).
\end{proof}

\begin{lemma}[Localization of an injective module]\label{res:L04}
For \(f\in R\) the map \(Q\to Q_f\) of an injective module is surjective
with kernel \(\Gamma_{(f)}(Q)\).
\end{lemma}
\begin{proof}
The kernel follows by the fraction relation. For \(e/f^a\in Q_f\) choose
\(b\) after the ascending sequence \(\operatorname{Ann}(f^j)\) has
stabilized. Then \(\operatorname{Ann}(f^{a+b})=\operatorname{Ann}(f^b)\).
The rule \(f^{a+b}r\mapsto rf^b e\) defines a map on the principal ideal
\((f^{a+b})\). Extend it to \(R\), and call its value at \(1\) \(q\).
In \(Q_f\) the equality \(f^{a+b}q=f^b e\) gives \(q=e/f^a\).
This argument also allows \(f=0\) or a unit.
\end{proof}

\section{The augmented \v{C}ech complex}

\begin{theorem}[Comparison with the augmented \v{C}ech complex]\label{res:L05}
For \(I=(f_1,\ldots,f_r)\) set
\[
 C_I^\bullet=\bigotimes_{j=1}^r [R\longrightarrow R_{f_j}],
\]
with each factor in degrees \(0,1\) and the empty tensor product \(R[0]\).
It computes \(R\Gamma_I\) naturally on bounded-below complexes by tensor
product. In particular \(H_I^q(M)=H^q(C_I^\bullet\otimes_R M)\).
\end{theorem}
\begin{proof}
The degree \(q\) term is the sum over \(q\)-element subsets \(J\) of the
localization at \(\prod_{j\in J}f_j\). Inserting index \(a\notin J\)
has sign \((-1)^{|\{j\in J:j<a\}|}\). Each pair of insertions in opposite
order has opposite sign, so this differential squares to zero.

For injective \(Q\), the inclusion
\(\Gamma_{(f)}(Q)[0]\to [Q\to Q_f]\) is a quasi-isomorphism by
Result~\ref{res:L04}. Its source is injective by Result~\ref{res:L03}.
Tensoring a quasi-isomorphism with a bounded complex of flat modules
preserves quasi-isomorphisms. Indeed, filter that bounded complex by
its degree tails and apply flatness to the single-term quotients
and the long exact sequence to their successive extensions.
Inducting over the factors now gives the natural quasi-isomorphism
\(\Gamma_I(Q)[0]\to C_I^\bullet\otimes Q\). Iteration of the principal
torsion functors gives \(\Gamma_I\) because individual annihilation
powers for finitely many \(f_j\) give a common power of \(I\).

Take a bounded-below injective replacement \(K\to J^\bullet\) as in
\cite[R01, R03 and R06]{homological}. The map
\(C_I^\bullet\otimes K\to C_I^\bullet\otimes J^\bullet\) is a
quasi-isomorphism by bounded flatness. The inclusion
\(\Gamma_I(J^\bullet)\to\operatorname{Tot}(C_I^\bullet\otimes J^\bullet)\)
is also a quasi-isomorphism because each augmented row is acyclic by the
injective case, and the bounded horizontal degree and lower bound
on \(J^\bullet\) give finite filtrations in each total degree.
This is the augmented totalization comparison of
\cite[B01--B04]{homological}. Both arrows are natural on replacements,
and replacement comparisons are unique in the derived category.
\end{proof}

\begin{corollary}[Generators and the cohomological bound]\label{res:L06}
The calculation is independent of the generators and depends only on
the radical of \(I\). If \(\sqrt I=\sqrt{(g_1,\ldots,g_s)}\), then
\(H_I^q(M)=0\) for \(q>s\). The zero and unit ideals have the derived
functors given by Result~\ref{res:L01}, with higher degrees zero.
\end{corollary}
\begin{proof}
Use the common functor of Result~\ref{res:L01} and the natural
comparisons of Result~\ref{res:L05}. The complex for the \(g_j\) has
length \(s\). Both extreme functors are exact.
\end{proof}

\begin{proposition}[The Koszul tower]\label{res:L07}
The complex \(C_I^\bullet\otimes M\) is the filtered colimit of the
cochain Koszul complexes on \(f_1^n,\ldots,f_r^n\) with coefficients in
\(M\). In the component \(J\) the transition to stage \(n+1\) is
multiplication by \(\prod_{j\in J}f_j\). In particular
\[
 H_I^r(M)\simeq
 \varinjlim_n M/(f_1^n,\ldots,f_r^n)M
\]
with transition multiplication by \(f_1\cdots f_r\).
\end{proposition}
\begin{proof}
The cochain differential in the component inserting \(a\) multiplies
by \(f_a^n\) with the sign of Result~\ref{res:L05}
(\cite[RK-03 and RK-10]{core}). The proposed maps commute with this
differential since \(f_a^{n+1}\prod_{j\in J}f_j
=f_a^n\prod_{j\in J\cup\{a\}}f_j\). A coefficient \(u\) at stage \(n\)
corresponds to \(u/(\prod_{j\in J}f_j)^n\). This identifies the
termwise colimit and the induced localization maps, including the
degree-zero component. The top cohomology at each stage is the stated
quotient. Apply filtered exactness. The ideals of powers of the
generators and the ordinary \(I\)-powers are cofinal by the same
pigeonhole argument as in Result~\ref{res:L01}.
\end{proof}

\section{Base change and the open complement}

\begin{proposition}[Change of rings and support]\label{res:L08}
Let \(R\to B\) be a homomorphism of Noetherian rings. For a \(B\)-module
\(N\), \(H_I^q(N)\simeq H_{IB}^q(N)\) compatibly with the \(B\)-action.
If \(B\) is \(R\)-flat then, for an \(R\)-module \(M\),
\[
 B\otimes_R H_I^q(M)\simeq H_{IB}^q(B\otimes_RM).
\]
In particular \(\operatorname{Supp} H_I^q(M)
\subseteq V(I)\cap\operatorname{Supp}M\).
\end{proposition}
\begin{proof}
Tensor the displayed localization terms of \(C_I^\bullet\) with \(B\).
Each is the localization of \(B\) at the images of the same generators.
Associativity of tensor product identifies the two complexes for \(N\).
For a flat map the tensor also takes the kernels and images in the
complex for \(M\) to the corresponding kernels and images after base
change. Localize at a prime to see that the cohomology is zero if \(M\)
localizes to zero or if \(I\) becomes the unit ideal.
\end{proof}

\begin{proposition}[Passage to the completion]\label{res:L09}
Let \(R\) be local and \(T\) be \(\mathfrak m\)-power torsion.
The \(R\)-action on \(T\) extends uniquely to \(\widehat R\), and
\(\widehat R\otimes_R T\to T\) is an isomorphism. For any \(R\)-module
\(M\), \(H^q_{\mathfrak m}(M)\) corresponds to
\(H^q_{\mathfrak m\widehat R}(\widehat R\otimes_RM)\).
\end{proposition}
\begin{proof}
On \(T[n]=(0:_T\mathfrak m^n)\), use
\(\widehat R/\mathfrak m^n\widehat R=R/\mathfrak m^n\)
from \cite[ALG-05-C04-S01]{p103}. The actions agree on overlaps,
and every extended action on an element killed by \(\mathfrak m^n\)
factors through that quotient. Tensor commutes with the ascending
union, giving the first isomorphism. Completion is flat by the same
prerequisite. Apply Result~\ref{res:L08} and then power torsion from
Result~\ref{res:L02}.
\end{proof}

\begin{proposition}[Cohomology of the complement]\label{res:L10}
Put \(U=\operatorname{Spec} R\setminus V(I)\).
There is a natural exact sequence
\[
 0\to H_I^0(M)\to M\to\Gamma(U,\widetilde M)\to H_I^1(M)\to0
\]
and \(H^q(U,\widetilde M)\simeq H_I^{q+1}(M)\) for \(q\geq1\).
\end{proposition}
\begin{proof}
Every open in a Noetherian affine spectrum has a finite distinguished
cover. Sections of an associated module sheaf on a distinguished open
are the corresponding localization. The sheaf axiom identifies sections
over the cover with the degree-zero cocycles of its unaugmented complex.
For an injective module \(Q\), Result~\ref{res:L05} shows that restriction
from \(Q\) to the sections on any open is surjective. This implies
that \(\widetilde Q\) is flasque, since an extension from a smaller
open can be restricted to any larger one. An injective module
resolution therefore gives a flasque resolution of \(\widetilde M\).
The sheafification of the module sequence is exact by exactness of
localization on stalks. By \cite[H01]{homological}, this gives
vanishing of positive sheaf cohomology on the affine spectrum and
on each distinguished affine intersection, applying the argument
over its localized Noetherian ring.

Now apply \cite[H02--H04]{homological} to the cover by \(D(f_j)\).
Its unaugmented complex is the complex of Result~\ref{res:L05}
in degrees \(\geq1\), shifted into degrees \(\geq0\), with the
usual sign comparison. Inserting \(M\) in degree zero gives precisely
the displayed low-degree sequence and the higher degree shift.
At complex level, if \(F^\bullet\) is the flasque replacement above,
the compatible \v{C}ech comparisons on \(X=\operatorname{Spec}R\) and
on \(U\) identify this augmented complex with the shifted cone of
\(\Gamma(X,F^\bullet)\to\Gamma(U,F^\bullet)\).
Indeed the augmentation of the double \v{C}ech complex on each flasque
term is exact, and the termwise restriction of \(M\to F^\bullet\)
is a resolution on every affine intersection with acyclic terms.
Taking the two comparisons in the total complex commutes with
the augmentation from \(X\). A homomorphism of modules lifts to a
map of injective resolutions, unique up to homotopy, so these
comparisons are natural on the derived complexes.
\end{proof}

\begin{corollary}[Mayer--Vietoris for supports]\label{res:L11}
For two ideals \(I,J\) there is a natural long exact sequence
\[
 \cdots\to H_{I+J}^q(M)\to H_I^q(M)\oplus H_J^q(M)
 \to H_{I\cap J}^q(M)\to H_{I+J}^{q+1}(M)\to\cdots .
\]
\end{corollary}
\begin{proof}
Take a flasque resolution \(F^\bullet\) of \(\widetilde M\) as in
Result~\ref{res:L10}. Write \(U=D(I)\), \(W=D(J)\) and \(X=\operatorname{Spec}R\).
The complexes for sections on \(U\cup W,U,W,U\cap W\) form the short
exact Mayer--Vietoris sequence of \cite[H05]{homological}, with difference
of restrictions in the last map. Augment the first and last terms by
\(\Gamma(X,F^\bullet)\) and the middle term by its two copies, using the
diagonal and the difference, respectively. Taking shifted mapping cones
of the augmented restrictions gives a short exact sequence of complexes.
Each cone computes the corresponding cohomology with support by the
complex comparison in Result~\ref{res:L10}. Here \(U\cup W=D(I+J)\)
and \(U\cap W=D(IJ)=D(I\cap J)\), which fixes the order of the terms.
\end{proof}
